\cotefile{157-2}
\scanpage{291}
\mspage{32}
\uline{Proposition} 7 Soit $f\colon I_{\unsure{\tau|t}}\longrightarrow I_{\tau'}$, un morphisme
d'intervalles \inserted{orientés.} \deleted{\unsure{Soit} \unsure{$f$|$s$} \unsure{une classe}} \deleted{Alors $f$ transforme sommet critique
en sommet critique \unsure{i.e.}}
\[
\deleted{f(I_{\tau\,\unsure{\mathrm{red}|\mathrm{r\acute ed}}})\supset(I_{\tau'})_{\mathrm{red}}}
\]

Soit
\[
 f\colon I_\tau\longrightarrow I_{\tau'}
\]
un morphisme d'intervalles orientés. \deleted{Les}
Considérons les sous-catégories pleines, \unsure{isom.}
à des $\unsure{\mathbb{\Delta}|\Delta}_i^+$ ou $\mathbb{\Delta}_i^-$ comme int. orientés,
\deleted{\unsure{qui}} définies par les intervalles critiques de
$I_\tau$. Ils forment une suite \inserted{$\Sigma_\tau$} de sous catégories
de $I_\tau$
\[
 \Sigma_\tau=(\deleted{\unsure{I_{\tau}}}\mathbb{\Delta}_{l_1}^{\varepsilon_1}(\tau)\unsure{\,|,}
 \mathbb{\Delta}_{l_2}^{\varepsilon_2}(\tau)\,\ldots,\mathbb{\Delta}_{l_r}^{\varepsilon_r}(\tau))
\]
dont $I_\tau$ est la somme amalgamée (on reviendra
sur ce point \unsure{plus}), $r$ étant la longueur utile de
$\tau$. On les appelle les \uline{sous}-\deleted{catégories}\uline{intervalles
critiques} de $I_\tau$. Ainsi $f$ se décompose en
une suite de morphismes
\[
 f_i\deleted{\unsure{I_\tau}}\colon\mathbb{\Delta}_{l_i}^{\varepsilon_i}\deleted{\unsure{\in}}
 \longrightarrow I_{\tau'},\qquad \deleted{\unsure{\forall}}i\in[1,r],
\]
tels que a) pour $i\in[1,r-1]$, le but de $f_i$
soit égal à la source de $f_{i+1}$. \deleted{\unsure{Et que}}\deleted{\unsure{soit d'entre}} De
plus, b) $f_1$ transforme source en source, et $f_r$ transforme
but en but. \deleted{Ces conditions} Inversement, quand
on se donne \unsure{des} $f_i$ satisfaisant ces deux conditions,
ils définissent un morphisme d'intervalles $f$
unique. (Et si la condition b) n'est pas \deleted{\unsure{sat}}
exigée, on trouve une description \unsure{des|de} morphismes
de catégories $I_\tau\longrightarrow I_{\tau'}$.)
% archive sheet number: 290
